\documentclass{ximera}

\input{../../../preamble.tex}


\definecolor{fvdnavy}{HTML}{071D43}
\definecolor{fvdcyan}{HTML}{20BFC6}
\definecolor{fvdcyanlight}{HTML}{EFFCFB}
\definecolor{fvdmagenta}{HTML}{F10D69}
\definecolor{fvdmagentalight}{HTML}{FFF1F7}
\definecolor{fvdtext}{HTML}{163044}





%=================================================
% Pedagogische omgevingen
% PDF: tcolorbox
% HTML: learning-box
%=================================================

\newenvironment{voorkennis}
{%
    \ifdefined\HCode
        \HCode{%
<section class="learning-box learning-box--prior">
<div class="learning-box__header">
<span class="learning-box__icon" aria-hidden="true"></span>
<span class="learning-box__title">Voorkennis</span>
</div>
<div class="learning-box__content">
        }%
        \begin{itemize}
    \else
       \begin{tcolorbox}[
    enhanced,
    breakable,
    colback=fvdcyanlight,
    colframe=fvdcyan,
    coltext=fvdtext,
    boxrule=0.5pt,
    leftrule=4pt,
    arc=4mm,
    outer arc=4mm,
    left=7mm,
    right=6mm,
    top=5mm,
    bottom=5mm,
    before skip=6mm,
    after skip=6mm,
    title=Voorkennis,
    coltitle=fvdcyan,
    fonttitle=\bfseries\Large,
    attach title to upper={\par\smallskip},
    boxed title style={
        empty,
        left=0mm,
        right=0mm,
        top=0mm,
        bottom=0mm
    },
    overlay={
        \draw[
            fvdcyan,
            line width=0.5pt
        ]
        ([xshift=42mm,yshift=-13mm]frame.north west)
        --
        ([xshift=-7mm,yshift=-13mm]frame.north east);
    }
]
        \begin{itemize}
    \fi
}
{%
    \end{itemize}
    \ifdefined\HCode
        \HCode{%
</div>
</section>
        }%
    \else
        \end{tcolorbox}
    \fi
}


\newenvironment{leerdoelen}
{%
    \ifdefined\HCode
        \HCode{%
<section class="learning-box learning-box--goals">
<div class="learning-box__header">
<span class="learning-box__icon" aria-hidden="true"></span>
<span class="learning-box__title">Leerdoelen</span>
</div>
<div class="learning-box__content">
        }%
        \begin{itemize}
    \else
\begin{tcolorbox}[
    enhanced,
    breakable,
    colback=fvdmagentalight,
    colframe=fvdmagenta,
    coltext=fvdtext,
    boxrule=0.5pt,
    leftrule=4pt,
    arc=4mm,
    outer arc=4mm,
    left=7mm,
    right=6mm,
    top=5mm,
    bottom=5mm,
    before skip=6mm,
    after skip=6mm,
    title=Leerdoelen,
    coltitle=fvdmagenta,
    fonttitle=\bfseries\Large,
    attach title to upper={\par\smallskip},
    boxed title style={
        empty,
        left=0mm,
        right=0mm,
        top=0mm,
        bottom=0mm
    },
    overlay={
        \draw[
            fvdmagenta,
            line width=0.5pt
        ]
        ([xshift=42mm,yshift=-13mm]frame.north west)
        --
        ([xshift=-7mm,yshift=-13mm]frame.north east);
    }
]
        \begin{itemize}
    \fi
}
{%
    \end{itemize}
    \ifdefined\HCode
        \HCode{%
</div>
</section>
        }%
    \else
        \end{tcolorbox}
    \fi
}


\addPrintStyle{../../..}

\title{Kracht en beweging}
\author{Ingmar Herreman}

\begin{abstract}
In de vorige hoofdstukken beschreven we hoe voorwerpen bewegen.
We onderzochten positie, snelheid en versnelling bij rechtlijnige,
tweedimensionale en cirkelvormige bewegingen. Nu stellen we een
fundamentelere vraag: waarom verandert een beweging?

Het antwoord ligt bij krachten. We bouwen de dynamica op vanuit de drie
wetten van Newton en leren krachten als vectoren samenstellen. Daarbij
staat niet het herkennen van losse formules centraal, maar het analyseren
van een volledig systeem: welk voorwerp onderzoeken we, welke krachten
werken erop, wat is hun resultante en welke versnelling volgt daaruit?

We besteden bijzondere aandacht aan vrijlichaamsdiagrammen, normaalkracht,
spankracht en wrijving. We leren ook enkele hardnekkige misconcepties
doorzien, zoals het verwarren van een Newton-III-paar met twee krachten
die elkaar op één voorwerp opheffen.

Ten slotte keren we terug naar de cirkelbeweging. De centripetale
versnelling uit het vorige hoofdstuk moet volgens Newton II veroorzaakt
worden door een naar binnen gerichte resulterende kracht. Daarmee leggen
we de basis voor de latere studie van gravitatie en satellietbeweging.
\end{abstract}

\begin{document}

% FVD_CHAPTER_DATA_START
% Automatisch gegenereerd uit index.tex.
% Niet handmatig aanpassen.
\fvdchapterdata{1}{Beweging en kracht}{6}
% FVD_CHAPTER_DATA_END

\maketitle

% ============================================================
% 1 VOORKENNIS EN LEERDOELEN
% ============================================================

\begin{voorkennis}
  \item Je kan positie, snelheid en versnelling vectorieel beschrijven.
  \item Je kan vectoren optellen en ontbinden in componenten.
  \item Je kan een eenparige en een eenparig versnelde beweging analyseren.
  \item Je kent de zwaartekracht nabij het aardoppervlak.
  \item Je kent bij een cirkelbeweging de richting van de centripetale
        versnelling en de formule $a_c=v^2/r$.
  \item Je kan goniometrie gebruiken om vectorcomponenten te bepalen.
\end{voorkennis}

\begin{leerdoelen}
  \item Je kan de drie wetten van Newton formuleren en in samenhang gebruiken.
  \item Je kan voor een gekozen systeem alle relevante uitwendige krachten
        identificeren en in een vrijlichaamsdiagram voorstellen.
  \item Je kan krachten vectorieel samenstellen en ontbinden.
  \item Je kan het vectoriële verband
        $\sum\vec F=m\vec a$ analyseren en kwantificeren.
  \item Je kan uit een gegeven krachtensituatie de bewegingsverandering
        voorspellen en omgekeerd uit de beweging conclusies trekken over
        de resulterende kracht.
  \item Je kan zwaartekracht, normaalkracht, spankracht en wrijvingskracht
        correct herkennen en modelleren.
  \item Je kan statische en kinetische wrijving onderscheiden.
  \item Je kan Newton-III-paren onderscheiden van krachten die op hetzelfde
        systeem werken.
  \item Je kan dynamische problemen in één en twee dimensies systematisch
        mathematiseren.
  \item Je kan de wetten van Newton toepassen op eenparige cirkelbeweging.
\end{leerdoelen}


% ============================================================
% 2 OPENING
% ============================================================

\FVDsection{Waarom verandert een beweging?}

Een auto trekt op, een parachutist valt, een lift vertrekt, een doos
wordt over een vloer geduwd en een satelliet verandert voortdurend van
bewegingsrichting.

In elk van deze situaties verandert de snelheidsvector:

\[
\vec v.
\]

Uit de kinematica weten we dat dit betekent dat er een versnelling is:

\[
\vec a\neq\vec 0.
\]

Maar de kinematica vertelt ons niet \emph{waarom} die versnelling
ontstaat. Daarvoor hebben we de dynamica nodig.

\begin{definitie}
De \textbf{dynamica} bestudeert het verband tussen krachten en de
bewegingsverandering van systemen.
\end{definitie}

De drie wetten van Newton vormen het fundament van de klassieke
dynamica.


\begin{oefening}[Voorspel vóór je rekent]{\oefB\ESS}

Een ijshockeypuck glijdt naar rechts over een vrijwel wrijvingsloze
horizontale ijsvlakte. Nadat de speler de puck heeft losgelaten, werkt
er geen horizontale kracht van de stick meer op.

\begin{question}
Is er volgens jou een blijvende kracht naar rechts nodig om de puck met
constante snelheid te laten voortbewegen?
\end{question}

\begin{question}
Wat verwacht je dat er gebeurt wanneer de horizontale resulterende
kracht nul is?
\end{question}

\begin{question}
Welke krachten werken verticaal op de puck?
\end{question}

Bewaar je eerste redenering. We komen hierop terug bij Newton I.

\end{oefening}


% ============================================================
% 3 SYSTEEM EN KRACHT
% ============================================================

\FVDsection{Begin met het systeem}

Voor we krachten tekenen, moeten we bepalen \emph{welk systeem} we
onderzoeken.

Dat kan bijvoorbeeld zijn:

\begin{itemize}
  \item één boek;
  \item een auto;
  \item een blok op een helling;
  \item twee verbonden massa's samen;
  \item of één van die massa's afzonderlijk.
\end{itemize}

De keuze van het systeem bepaalt welke krachten uitwendig zijn en dus
in het krachtendiagram thuishoren.

\begin{definitie}
Een \textbf{kracht} is een vectoriële wisselwerking die de
bewegingsverandering van een systeem kan beïnvloeden.
\end{definitie}

Een kracht heeft dus:

\[
\boxed{\text{grootte, richting en zin}}
\]

en wordt voorgesteld door een vector $\vec F$.

De SI-eenheid is de newton:

\[
[F]=\mathrm N.
\]


\FVDsubsection{Enkele krachten die we vaak gebruiken}

In dit hoofdstuk ontmoeten we onder andere:

\[
\vec F_g
\qquad \text{zwaartekracht},
\]

\[
\vec F_N
\qquad \text{normaalkracht},
\]

\[
\vec F_T
\qquad \text{spankracht},
\]

en

\[
\vec F_w
\qquad \text{wrijvingskracht}.
\]

Deze krachten zijn geen formules die automatisch moeten worden ingevuld.
We moeten telkens eerst onderzoeken of de betrokken wisselwerking
werkelijk aanwezig is.


% ============================================================
% 4 RESULTANTE
% ============================================================

\FVDsection{De resulterende kracht}

Op een systeem kunnen meerdere krachten tegelijk werken.

De \textbf{resulterende kracht} is hun vectoriële som:

\[
\boxed{
\vec F_{\mathrm{res}}
=
\sum_i\vec F_i.
}
\]

In twee dimensies analyseren we componentgewijs:

\[
\boxed{
\sum F_x=F_{\mathrm{res},x}
}
\]

en

\[
\boxed{
\sum F_y=F_{\mathrm{res},y}.
}
\]

Een nulresultante betekent:

\[
\sum\vec F=\vec 0.
\]

Dat betekent niet dat er geen krachten werken. Het betekent dat de
vectoriële som van alle uitwendige krachten nul is.


\begin{oefening}[Krachten samenstellen]{\oefB\ESS}

Op een voorwerp werken horizontaal drie krachten:

\[
F_1=18\,\mathrm N\quad\text{naar rechts},
\]

\[
F_2=7,0\,\mathrm N\quad\text{naar links},
\]

\[
F_3=4,0\,\mathrm N\quad\text{naar links}.
\]

\begin{question}
Kies rechts als positieve richting en schrijf de drie krachten met teken.
\end{question}

\begin{question}
Bepaal de resulterende kracht.
\end{question}

\begin{question}
Wat kan je al zeggen over de richting van de versnelling?
\end{question}

\end{oefening}


% ============================================================
% 5 NEWTON I
% ============================================================

\FVDsection{De eerste wet van Newton: traagheid}

Een hardnekkige intuïtie is dat een voorwerp voortdurend een kracht in
de bewegingsrichting nodig heeft om te blijven bewegen.

Newton I zegt iets anders.

\begin{eigenschap}[Eerste wet van Newton]
Wanneer de resulterende uitwendige kracht op een systeem nul is, blijft
de snelheid van het systeem constant:

\[
\boxed{
\sum\vec F=\vec 0
\quad\Longrightarrow\quad
\vec v=\text{constant}.
}
\]

Een systeem in rust blijft in rust; een bewegend systeem blijft
rechtlijnig met constante snelheid bewegen.
\end{eigenschap}

Dit noemen we het \textbf{traagheidsbeginsel}.

Merk op:

\[
\vec v=\text{constant}
\]

kan ook betekenen dat $\vec v=\vec 0$.


\FVDsubsection{Evenwicht is niet hetzelfde als rust}

Als

\[
\sum\vec F=\vec 0,
\]

dan is

\[
\vec a=\vec 0.
\]

Maar het voorwerp hoeft daarom niet stil te staan.

Het kan:

\begin{itemize}
  \item in rust zijn;
  \item of met constante snelheid langs een rechte bewegen.
\end{itemize}


\begin{oefening}[Kracht nodig om te blijven bewegen?]{\oefU\ESS}

Een ruimtesonde bevindt zich ver van grote hemellichamen. Haar motoren
worden uitgeschakeld terwijl ze een snelheid $\vec v_0$ heeft.

Verwaarloos alle uitwendige krachten.

\begin{question}
Beschrijf de beweging na het uitschakelen van de motoren.
\end{question}

\begin{question}
Is er een kracht in de bewegingsrichting nodig om $\vec v_0$ te
behouden?
\end{question}

\begin{question}
Welke misconceptie wordt door deze situatie zichtbaar?
\end{question}

\end{oefening}


% ============================================================
% 6 NEWTON II
% ============================================================

\FVDsection{De tweede wet van Newton}

Newton II verbindt de resulterende kracht met de versnelling.

\begin{eigenschap}[Tweede wet van Newton]

Voor een systeem met constante massa geldt:

\[
\boxed{
\sum\vec F=m\vec a.
}
\]

\end{eigenschap}

Dit is een \textbf{vectorvergelijking}. Ze bevat tegelijk informatie
over grootte én richting.

Omdat $m>0$ geldt:

\[
\boxed{
\vec a\ \text{heeft dezelfde richting en zin als}\ 
\sum\vec F.
}
\]

Niet noodzakelijk als de snelheid.


\FVDsubsection{Componentvorm}

In een cartesisch assenstelsel:

\[
\boxed{
\sum F_x=ma_x
}
\]

en

\[
\boxed{
\sum F_y=ma_y.
}
\]

Dit is één van de belangrijkste analysetechnieken van de mechanica.


\FVDsubsection{Wat vertelt Newton II kwalitatief?}

Bij dezelfde massa:

\[
a\sim F_{\mathrm{res}}.
\]

Een tweemaal zo grote resulterende kracht geeft dus een tweemaal zo
grote versnelling.

Bij dezelfde resulterende kracht:

\[
a\sim\frac1m.
\]

Een tweemaal zo grote massa geeft dus een half zo grote versnelling.


\begin{oefening}[Van verband naar voorspelling]{\oefU\ESS}

Een karretje met massa $m$ ondervindt een constante horizontale
resulterende kracht $F$ en krijgt versnelling $a$.

Redeneer zonder concrete getallen.

\begin{question}
Wat gebeurt er met $a$ als $F$ verdrievoudigt en $m$ gelijk blijft?
\end{question}

\begin{question}
Wat gebeurt er met $a$ als $m$ verdubbelt en $F$ gelijk blijft?
\end{question}

\begin{question}
Wat gebeurt er met $a$ als zowel $F$ als $m$ verdubbelen?
\end{question}

\begin{question}
Waarom moet je bij zulke redeneringen altijd aangeven welke grootheid
constant blijft?
\end{question}

\end{oefening}


\FVDsubsection{Constante kracht en EVRB}

Wanneer de massa constant is en de resulterende kracht constant blijft:

\[
\sum\vec F=\text{constant},
\]

dan volgt uit Newton II:

\[
\vec a=\text{constant}.
\]

In één dimensie herkennen we dan de eenparig versnelde rechtlijnige
beweging uit de vorige hoofdstukken.

Dynamica en kinematica sluiten dus rechtstreeks op elkaar aan:

\[
\boxed{
\sum F=\text{constant}
\Longrightarrow
a=\text{constant}
\Longrightarrow
\text{EVRB}.
}
\]


\begin{oefening}[Van kracht naar bewegingsfunctie]{\oefU\ESS}

Een karretje met massa

\[
m=2,0\,\mathrm{kg}
\]

start vanuit rust op $x_0=0$. Vanaf $t=0$ werkt een constante
resulterende horizontale kracht van $6,0\,\mathrm N$ naar rechts.

\begin{question}
Bepaal de versnelling.
\end{question}

\begin{question}
Stel $v_x(t)$ op.
\end{question}

\begin{question}
Stel $x(t)$ op.
\end{question}

\begin{question}
Bepaal de positie en snelheid na $4,0\,\mathrm s$.
\end{question}

\begin{question}
Leg uit hoe in deze oefening Newton II en de kinematica samenwerken.
\end{question}

\end{oefening}


% ============================================================
% 7 VRIJLICHAAMSDIAGRAM
% ============================================================

\FVDsection{Het vrijlichaamsdiagram}

Een betrouwbaar Newtonprobleem begint niet met een formule, maar met een
model van de krachten.

Een \textbf{vrijlichaamsdiagram} (VLD) bevat alleen het gekozen systeem
en de uitwendige krachten die daarop werken.

Een goede werkwijze is:

\begin{enumerate}
  \item Kies het systeem.
  \item Isoleer het systeem denkbeeldig van zijn omgeving.
  \item Zoek alle wisselwerkingen met de omgeving.
  \item Teken voor elke wisselwerking één krachtvector.
  \item Benoem de krachten.
  \item Kies pas daarna geschikte assen.
\end{enumerate}

\begin{opmerking}
Teken in een VLD geen snelheid of versnelling alsof het krachten zijn.
$\vec v$ en $\vec a$ beschrijven de beweging; $\vec F$ beschrijft een
wisselwerking.
\end{opmerking}


\begin{oefening}[Welke krachten horen erbij?]{\oefB\ESS}

Een boek ligt stil op een horizontale tafel.

\begin{question}
Kies het boek als systeem.
\end{question}

\begin{question}
Teken het vrijlichaamsdiagram.
\end{question}

\begin{question}
Welke kracht oefent de aarde op het boek uit?
\end{question}

\begin{question}
Welke kracht oefent de tafel op het boek uit?
\end{question}

\begin{question}
Wat is de verticale resultante?
\end{question}

\end{oefening}


% ============================================================
% 8 ZWAARTEKRACHT
% ============================================================

\FVDsection{Zwaartekracht nabij het aardoppervlak}

Nabij het aardoppervlak modelleren we de zwaartekracht op een massa $m$
als:

\[
\boxed{
\vec F_g=m\vec g.
}
\]

De grootte is:

\[
\boxed{
F_g=mg.
}
\]

De richting is verticaal en de zin is naar het middelpunt van de aarde.

In een assenstelsel met omhoog positief:

\[
F_{g,y}=-mg.
\]

\begin{opmerking}
Massa en gewicht zijn niet hetzelfde. Massa $m$ is een eigenschap van
het systeem en wordt uitgedrukt in kilogram. De zwaartekracht $F_g$ is
een kracht en wordt uitgedrukt in newton.
\end{opmerking}


\begin{oefening}[Vrije val opnieuw bekeken]{\oefU\ESS}

Een voorwerp valt in vacuüm nabij het aardoppervlak.

\begin{question}
Welke kracht werkt op het voorwerp als we andere invloeden
verwaarlozen?
\end{question}

\begin{question}
Gebruik Newton II om aan te tonen dat de versnelling onafhankelijk is
van de massa.
\end{question}

\begin{question}
Verklaar waarom een zwaarder voorwerp in dit model niet sneller valt
dan een lichter voorwerp.
\end{question}

\end{oefening}


% ============================================================
% 9 NORMAALKRACHT
% ============================================================

\FVDsection{De normaalkracht}

Wanneer twee oppervlakken contact maken, kan het oppervlak een
contactkracht uitoefenen loodrecht op het contactvlak.

Die loodrechte component noemen we de \textbf{normaalkracht}:

\[
\vec F_N.
\]

``Normaal'' betekent hier: loodrecht.


\FVDsubsection{De normaalkracht is niet altijd $mg$}

Voor een boek in rust op een horizontale tafel, zonder andere verticale
krachten:

\[
\sum F_y=0
\]

geeft:

\[
F_N-F_g=0.
\]

Dus in deze specifieke situatie:

\[
F_N=mg.
\]

Maar dit is geen algemene natuurwet.

Wanneer bijvoorbeeld extra verticaal wordt geduwd of getrokken, of
wanneer het systeem verticaal versnelt, kan $F_N$ verschillen van
$mg$.


\begin{oefening}[Wanneer is $F_N=mg$?]{\oefU\ESS}

Een kist met massa $20\,\mathrm{kg}$ staat op een horizontale vloer.

Onderzoek drie situaties.

\begin{question}
Er werken geen andere verticale krachten en de kist versnelt niet
verticaal. Bepaal $F_N$.
\end{question}

\begin{question}
Iemand duwt bovendien verticaal naar beneden met $50\,\mathrm N$.
Bepaal $F_N$ zolang de kist niet verticaal versnelt.
\end{question}

\begin{question}
Iemand trekt verticaal omhoog met $50\,\mathrm N$, maar de kist blijft
op de vloer. Bepaal $F_N$.
\end{question}

\begin{question}
Formuleer waarom de regel ``$F_N=mg$'' gevaarlijk is.
\end{question}

\end{oefening}


% ============================================================
% 10 SPANKRACHT
% ============================================================

\FVDsection{Spankracht}

Een gespannen touw, kabel of koord kan een trekkracht overbrengen.

We noemen die kracht de \textbf{spankracht}:

\[
\vec F_T.
\]

Een ideaal massaloos touw met ideale katrollen wordt vaak gemodelleerd
met dezelfde spankracht langs het touw. In echte touwen en niet-ideale
katrollen kan dat model beperkingen hebben.

\begin{oefening}[Hangende massa]{\oefB}

Een lamp met massa $12\,\mathrm{kg}$ hangt stil aan één verticale kabel.

\begin{question}
Teken het VLD van de lamp.
\end{question}

\begin{question}
Bepaal de spankracht in de kabel.
\end{question}

\begin{question}
Welke twee krachten vormen hier een krachtenevenwicht op de lamp?
\end{question}

\end{oefening}


% ============================================================
% 11 NEWTON III
% ============================================================

\FVDsection{De derde wet van Newton}

Krachten ontstaan door wisselwerkingen tussen systemen.

Als systeem A een kracht uitoefent op systeem B, dan oefent B
tegelijk een kracht uit op A.

\begin{eigenschap}[Derde wet van Newton]

\[
\boxed{
\vec F_{A\rightarrow B}
=
-\vec F_{B\rightarrow A}.
}
\]

De twee krachten hebben dezelfde grootte, dezelfde werklijn en
tegengestelde zin, maar werken op \textbf{verschillende systemen}.
\end{eigenschap}


\FVDsubsection{Newton III is geen krachtenevenwicht}

Een boek ligt op een tafel.

Op het boek werken onder andere:

\[
\vec F_{g,\mathrm{aarde}\rightarrow\mathrm{boek}}
\]

en

\[
\vec F_{N,\mathrm{tafel}\rightarrow\mathrm{boek}}.
\]

Deze twee krachten kunnen elkaar opheffen, maar ze vormen geen
Newton-III-paar: ze werken allebei op het boek.

Het Newton-III-paar bij de normaalkracht is:

\[
\vec F_{\mathrm{tafel}\rightarrow\mathrm{boek}}
=
-\vec F_{\mathrm{boek}\rightarrow\mathrm{tafel}}.
\]

En bij de gravitatie-interactie:

\[
\vec F_{\mathrm{aarde}\rightarrow\mathrm{boek}}
=
-\vec F_{\mathrm{boek}\rightarrow\mathrm{aarde}}.
\]


\begin{oefening}[Actie-reactie of evenwicht?]{\oefU\ESS}

Een persoon staat stil op de vloer.

\begin{question}
Teken het VLD van de persoon.
\end{question}

\begin{question}
Zijn de zwaartekracht op de persoon en de normaalkracht van de vloer
op de persoon een Newton-III-paar? Verklaar.
\end{question}

\begin{question}
Geef het Newton-III-paar van de normaalkracht.
\end{question}

\begin{question}
Geef het Newton-III-paar van de zwaartekracht op de persoon.
\end{question}

\end{oefening}


\begin{oefening}[Paard en kar]{\oefU\ESS}

Een leerling redeneert:

\begin{quote}
``Het paard trekt aan de kar en de kar trekt volgens Newton III even
hard terug aan het paard. Die krachten heffen elkaar dus op. De kar kan
nooit versnellen.''
\end{quote}

\begin{question}
Wat is fout aan deze redenering?
\end{question}

\begin{question}
Op welk systeem werkt de kracht van het paard op de kar?
\end{question}

\begin{question}
Op welk systeem werkt de reactiekracht van de kar op het paard?
\end{question}

\begin{question}
Waarom mogen deze twee krachten niet in één som
$\sum\vec F$ voor één systeem worden opgeteld?
\end{question}

\end{oefening}


% ============================================================
% 12 WRIJVING
% ============================================================

\FVDsection{Wrijvingskracht}

Wanneer twee oppervlakken contact maken, kan naast de normaalkracht ook
een kracht evenwijdig aan het contactvlak optreden: de
\textbf{wrijvingskracht}.

Wrijving verzet zich tegen relatieve glijbeweging of tegen de neiging
tot glijden tussen de contactoppervlakken.


\FVDsubsection{Statische wrijving}

Als de oppervlakken niet ten opzichte van elkaar glijden, spreken we
van \textbf{statische wrijving}.

De statische wrijvingskracht past zich aan de situatie aan tot een
maximale waarde:

\[
\boxed{
0\leq F_{w,s}\leq \mu_sF_N.
}
\]

De maximale statische wrijving is:

\[
\boxed{
F_{w,s,\max}=\mu_sF_N.
}
\]

Belangrijk:

\[
F_{w,s}=\mu_sF_N
\]

geldt dus alleen wanneer de statische wrijving haar maximale waarde
bereikt, bijvoorbeeld vlak vóór het glijden begint.


\FVDsubsection{Kinetische wrijving}

Wanneer oppervlakken over elkaar glijden, gebruiken we in het
eenvoudige model:

\[
\boxed{
F_{w,k}=\mu_kF_N.
}
\]

De kracht is tegengesteld aan de relatieve glijbeweging.

Vaak geldt:

\[
\mu_s>\mu_k,
\]

maar dit is een materiaaleigenschap en geen universele wiskundige
noodzaak.


\begin{oefening}[De wrijving kiest niet automatisch $\mu_sF_N$]{\oefU\ESS}

Een kist staat op een horizontale vloer. De maximale statische
wrijvingskracht bedraagt $80\,\mathrm N$.

\begin{question}
Je duwt horizontaal met $30\,\mathrm N$ en de kist blijft stil.
Hoe groot is de statische wrijvingskracht?
\end{question}

\begin{question}
Je duwt met $60\,\mathrm N$ en de kist blijft stil.
Hoe groot is de statische wrijvingskracht?
\end{question}

\begin{question}
Waarom is in beide gevallen $F_{w,s}\neq80\,\mathrm N$?
\end{question}

\begin{question}
Wat betekent de waarde $80\,\mathrm N$ fysisch?
\end{question}

\end{oefening}


\begin{oefening}[Van rust naar glijden]{\oefU\ESS}

Een kist met massa $25\,\mathrm{kg}$ staat op een horizontale vloer.
Neem:

\[
\mu_s=0,40,
\qquad
\mu_k=0,30.
\]

Er wordt horizontaal aan de kist getrokken.

\begin{question}
Bepaal $F_N$.
\end{question}

\begin{question}
Bepaal de maximale statische wrijvingskracht.
\end{question}

\begin{question}
Zal de kist bewegen bij een trekkracht van $80\,\mathrm N$?
Bepaal de werkelijke wrijvingskracht.
\end{question}

\begin{question}
Zal de kist bewegen bij een trekkracht van $120\,\mathrm N$?
\end{question}

\begin{question}
Bepaal in dat laatste geval de kinetische wrijvingskracht en de
versnelling nadat de kist begint te glijden.
\end{question}

\end{oefening}


% ============================================================
% 13 WRIJVING EN BEWEGING
% ============================================================

\FVDsection{Wrijving werkt niet altijd ``tegen de beweging''}

De eenvoudige zin ``wrijving werkt tegen de beweging'' kan misleidend
zijn.

Denk aan een auto die versnelt. De aangedreven banden duwen via hun
contact met de weg naar achteren op het wegdek. Het wegdek oefent via
statische wrijving een kracht naar voren uit op de banden.

Die wrijvingskracht kan dus dezelfde zin hebben als de beweging van de
auto.

De nauwkeurigere formulering is:

\[
\boxed{
\text{wrijving verzet zich tegen relatieve glijbeweging of glijneiging.}
}
\]


\begin{oefening}[Waarom kan een auto optrekken?]{\oefG}

Een auto trekt op een horizontale weg op zonder dat de aangedreven
banden slippen.

\begin{question}
Is de relevante wrijving tussen band en weg statisch of kinetisch?
\end{question}

\begin{question}
In welke zin duwt de aangedreven band op het wegdek?
\end{question}

\begin{question}
Gebruik Newton III om de zin van de kracht van het wegdek op de band
te bepalen.
\end{question}

\begin{question}
Leg uit waarom de uitspraak ``wrijving werkt altijd tegen de
bewegingsrichting'' hier tekortschiet.
\end{question}

\end{oefening}


% ============================================================
% 14 KRACHTEN ONTBINDEN
% ============================================================

\FVDsection{Krachten ontbinden}

In twee dimensies is het vaak nuttig een kracht in componenten te
ontbinden.

Als een kracht $\vec F$ een hoek $\alpha$ maakt met de positieve
$x$-as:

\[
\boxed{
F_x=F\cos\alpha
}
\]

en

\[
\boxed{
F_y=F\sin\alpha.
}
\]

Daarna passen we Newton II per richting toe:

\[
\boxed{
\sum F_x=ma_x,
\qquad
\sum F_y=ma_y.
}
\]


\FVDsubsection{Slimme assen kiezen}

De assen hoeven niet altijd horizontaal en verticaal te liggen.

Bij een hellend vlak is het vaak efficiënter om:

\begin{itemize}
  \item de $x$-as evenwijdig aan het vlak;
  \item de $y$-as loodrecht op het vlak
\end{itemize}

te kiezen.

Dan ligt de normaalkracht volledig langs één as.


% ============================================================
% 15 HELLEND VLAK
% ============================================================

\FVDsection{Het hellend vlak}

Een blok ligt op een helling met hoek $\alpha$.

De zwaartekracht wijst verticaal naar beneden:

\[
F_g=mg.
\]

Met assen evenwijdig en loodrecht op de helling ontbinden we:

\[
\boxed{
F_{g,\parallel}=mg\sin\alpha
}
\]

en

\[
\boxed{
F_{g,\perp}=mg\cos\alpha.
}
\]

Als het blok niet loodrecht op het vlak versnelt:

\[
\sum F_\perp=0.
\]

Zonder andere loodrechte krachten volgt:

\[
\boxed{
F_N=mg\cos\alpha.
}
\]

Opnieuw zien we dat $F_N$ niet automatisch gelijk is aan $mg$.


\begin{oefening}[Wrijvingsloos hellend vlak]{\oefU\ESS}

Een blok met massa $4,0\,\mathrm{kg}$ glijdt zonder wrijving langs een
helling met hoek $30^\circ$.

\begin{question}
Teken het VLD.
\end{question}

\begin{question}
Kies geschikte assen en ontbind $\vec F_g$.
\end{question}

\begin{question}
Bepaal $F_N$.
\end{question}

\begin{question}
Bepaal de resulterende kracht langs de helling.
\end{question}

\begin{question}
Bepaal de versnelling.
\end{question}

\begin{question}
Waarom verdwijnt de massa uit het eindresultaat voor $a$?
\end{question}

\end{oefening}


\begin{oefening}[Hellend vlak met wrijving]{\oefU\ESS}

Een blok glijdt naar beneden langs een helling met hoek $\alpha$.
De kinetische wrijvingscoëfficiënt is $\mu_k$.

\begin{question}
Teken een algemeen VLD.
\end{question}

\begin{question}
Toon aan dat, zolang er geen andere krachten werken,

\[
F_N=mg\cos\alpha.
\]

\end{question}

\begin{question}
Toon aan dat de kinetische wrijving de grootte

\[
F_{w,k}=\mu_kmg\cos\alpha
\]

heeft.
\end{question}

\begin{question}
Leid af dat de versnelling langs het vlak gegeven wordt door

\[
\boxed{
a=g(\sin\alpha-\mu_k\cos\alpha).
}
\]

\end{question}

\begin{question}
Wat valt je op aan de massa? Interpreteer.
\end{question}

\end{oefening}


% ============================================================
% 16 SYSTEMATISCHE PROBLEEMAANPAK
% ============================================================

\FVDsection{Newtonproblemen mathematiseren}

Bij complexere problemen is de grootste uitdaging meestal niet het
algebraïsche rekenwerk, maar het vertalen van de situatie naar een
correct fysisch model.

Gebruik daarom systematisch:

\begin{enumerate}
  \item \textbf{Lees en schets.} Wat gebeurt er fysisch?
  \item \textbf{Kies het systeem.} Over welk voorwerp of welke verzameling
        voorwerpen gaat $\sum\vec F=m\vec a$?
  \item \textbf{Maak een VLD.} Teken alleen de uitwendige krachten op dat
        systeem.
  \item \textbf{Kies assen.} Kies ze zo dat zoveel mogelijk krachten of
        de versnelling langs een as liggen.
  \item \textbf{Ontbind vectoren.}
  \item \textbf{Schrijf Newton II per richting:}
        \[
        \sum F_x=ma_x,\qquad \sum F_y=ma_y.
        \]
  \item \textbf{Voeg modelrelaties toe.} Bijvoorbeeld
        $F_g=mg$, $F_{w,k}=\mu_kF_N$ of een kinematische vergelijking.
  \item \textbf{Los het stelsel op.}
  \item \textbf{Controleer.} Kloppen teken, richting, eenheid en
        grootteorde?
  \item \textbf{Demathematiseer.} Wat betekent het resultaat in de
        oorspronkelijke situatie?
\end{enumerate}

\begin{opmerking}
Een negatieve uitkomst voor een component betekent niet automatisch dat
de berekening fout is. Ze kan aangeven dat de werkelijke zin tegengesteld
is aan de gekozen positieve richting.
\end{opmerking}


% ============================================================
% 17 VOORBEELD VOLLEDIGE ANALYSE
% ============================================================

\FVDsection{Voorbeeld: trekken aan een kist}

Een kist met massa $m$ wordt met een kracht $\vec F$ onder hoek
$\alpha$ boven de horizontale over een ruwe horizontale vloer getrokken.

Dit probleem toont waarom een VLD belangrijker is dan formules
memoriseren.

Verticaal is er geen versnelling:

\[
a_y=0.
\]

Dus:

\[
\sum F_y=0.
\]

Met omhoog positief:

\[
F_N+F\sin\alpha-mg=0.
\]

Daaruit:

\[
\boxed{
F_N=mg-F\sin\alpha.
}
\]

De trekkracht vermindert dus de normaalkracht.

Bij kinetische wrijving:

\[
F_{w,k}=\mu_kF_N
\]

en dus:

\[
F_{w,k}
=
\mu_k(mg-F\sin\alpha).
\]

Horizontaal:

\[
F\cos\alpha-F_{w,k}=ma_x.
\]

Daarmee kunnen we de versnelling bepalen.

Dit voorbeeld toont een belangrijk inzicht:

\[
\boxed{
F_w=\mu F_N
}
\]

kan niet correct worden gebruikt voordat $F_N$ uit de volledige
krachtensituatie is bepaald.


\begin{oefening}[Trekkracht onder een hoek]{\oefU\ESS}

Een kist met massa $30\,\mathrm{kg}$ wordt over een horizontale vloer
getrokken met een kracht van $150\,\mathrm N$ onder een hoek van
$25^\circ$ boven de horizontale. Neem $\mu_k=0,20$.

\begin{question}
Teken het VLD.
\end{question}

\begin{question}
Ontbind de trekkracht.
\end{question}

\begin{question}
Bepaal de normaalkracht.
\end{question}

\begin{question}
Bepaal de kinetische wrijvingskracht.
\end{question}

\begin{question}
Bepaal de horizontale versnelling.
\end{question}

\begin{question}
Vergelijk kwalitatief met dezelfde trekkracht die volledig horizontaal
zou werken. Waarom verandert de wrijving?
\end{question}

\end{oefening}


% ============================================================
% 18 TWEE SYSTEMEN
% ============================================================

\FVDsection{Meerdere systemen analyseren}

Soms zijn verschillende voorwerpen gekoppeld.

Dan kunnen we:

\begin{itemize}
  \item elk voorwerp afzonderlijk als systeem kiezen;
  \item of de volledige verzameling als één groter systeem beschouwen.
\end{itemize}

Beide keuzes kunnen nuttig zijn.


\begin{oefening}[Twee blokken]{\oefU\ESS}

Twee blokken met massa's $m_1$ en $m_2$ liggen op een wrijvingsloze
horizontale tafel en zijn verbonden door een massaloos touw. Op
$m_2$ werkt horizontaal een uitwendige kracht $F$.

\begin{question}
Beschouw eerst beide blokken samen als één systeem. Leid de versnelling
af.
\end{question}

\begin{question}
Beschouw daarna alleen blok $m_1$. Bepaal de spankracht in het touw.
\end{question}

\begin{question}
Toon dat

\[
a=\frac{F}{m_1+m_2}.
\]

\end{question}

\begin{question}
Toon dat

\[
F_T=\frac{m_1}{m_1+m_2}F.
\]

\end{question}

\begin{question}
Waarom verschijnt de interne spankracht niet wanneer beide blokken samen
als systeem worden gekozen?
\end{question}

\end{oefening}


% ============================================================
% 19 NEWTON I II III SAMEN
% ============================================================

\FVDsection{De drie wetten werken samen}

In echte analyses gebruik je de wetten van Newton zelden als drie
volledig losstaande regels.

Newton I helpt herkennen wat een nulresultante betekent.

Newton II verbindt de resultante met de versnelling:

\[
\sum\vec F=m\vec a.
\]

Newton III helpt wisselwerkingen tussen verschillende systemen correct
te identificeren.

Bij een auto die versnelt, bijvoorbeeld:

\begin{itemize}
  \item de banden oefenen een kracht uit op de weg;
  \item de weg oefent volgens Newton III een tegengestelde kracht uit op
        de banden;
  \item die kracht draagt bij aan de resulterende kracht op de auto;
  \item volgens Newton II veroorzaakt de resultante de versnelling.
\end{itemize}

Dit is precies waarom een beweging vaak alleen goed begrepen wordt door
de drie wetten samen te gebruiken.


\begin{oefening}[Liftanalyse]{\oefU\ESS}

Een persoon met massa $70\,\mathrm{kg}$ staat op een weegschaal in een
lift. De weegschaal meet de normaalkracht.

Analyseer telkens met omhoog positief.

\begin{question}
Bepaal de aflezing wanneer de lift stilstaat.
\end{question}

\begin{question}
Bepaal de aflezing wanneer de lift met
$1,5\,\mathrm{m/s^2}$ omhoog versnelt.
\end{question}

\begin{question}
Bepaal de aflezing wanneer de lift met
$1,5\,\mathrm{m/s^2}$ naar beneden versnelt.
\end{question}

\begin{question}
Wat leest de weegschaal af tijdens vrije val in het ideale model?
\end{question}

\begin{question}
Waarom verandert de massa van de persoon in geen van deze situaties?
\end{question}

\end{oefening}


% ============================================================
% 20 TERUG NAAR CIRKELBEWEGING
% ============================================================

\FVDsection{Newton en de cirkelbeweging}

In het vorige hoofdstuk vonden we voor een ECB:

\[
\boxed{
a_c=\frac{v^2}{r}.
}
\]

De versnelling wijst naar het middelpunt.

Newton II zegt dan dat ook de radiale resulterende kracht naar het
middelpunt moet wijzen:

\[
\sum F_r=ma_c.
\]

Dus:

\[
\boxed{
\sum F_r
=
m\frac{v^2}{r}.
}
\]

Dit is geen nieuwe krachtsoort. Het is Newton II toegepast in de
radiale richting.


\FVDsubsection{Geen aparte ``centripetale kracht''}

De naar binnen gerichte resultante kan geleverd worden door verschillende
echte krachten of combinaties daarvan.

Bijvoorbeeld:

\begin{itemize}
  \item spankracht bij een massa aan een touw;
  \item wrijving bij een auto in een vlakke bocht;
  \item normaalkracht in bepaalde gebogen banen;
  \item gravitatie bij een satelliet.
\end{itemize}

Daarom is de formulering

\[
F_c=\frac{mv^2}{r}
\]

gevaarlijk wanneer ze doet vermoeden dat $F_c$ naast alle andere
krachten in het VLD moet worden getekend.

Beter is:

\[
\boxed{
\sum F_r=\frac{mv^2}{r}.
}
\]


\begin{oefening}[Auto in een vlakke bocht]{\oefU\ESS}

Een auto met massa $1200\,\mathrm{kg}$ rijdt met constante
snelheidsgrootte door een vlakke cirkelvormige bocht met straal
$50\,\mathrm m$. De snelheid bedraagt $15\,\mathrm{m/s}$.

\begin{question}
Teken het VLD.
\end{question}

\begin{question}
Welke kracht kan in het eenvoudige model de horizontale radiale
resultante leveren?
\end{question}

\begin{question}
Bepaal de vereiste radiale resultante.
\end{question}

\begin{question}
Bepaal de minimale statische wrijvingscoëfficiënt om niet te slippen.
\end{question}

\begin{question}
Wat gebeurt er met de vereiste radiale kracht als de snelheid
verdubbelt?
\end{question}

\end{oefening}


\begin{oefening}[Massa aan een touw]{\oefU}

Een massa beweegt met constante snelheid in een horizontale cirkel.
Veronderstel in een vereenvoudigd model dat de spankracht de volledige
horizontale radiale resultante levert.

\begin{question}
Leid af:

\[
F_T=\frac{mv^2}{r}.
\]

\end{question}

\begin{question}
Hoe verandert $F_T$ als $v$ verdubbelt?
\end{question}

\begin{question}
Hoe verandert $F_T$ als $r$ verdubbelt terwijl $v$ constant blijft?
\end{question}

\begin{question}
Waarom moet je in een echte opstelling altijd eerst het VLD bekijken
voordat je de spankracht gelijkstelt aan $mv^2/r$?
\end{question}

\end{oefening}


% ============================================================
% 21 GRAVITATIE ALS VOORUITBLIK
% ============================================================

\FVDsection{Vooruitblik: waarom valt een satelliet niet naar beneden?}

Een satelliet in een cirkelbaan heeft een centripetale versnelling naar
het middelpunt van de aarde.

Volgens Newton II moet er dus een naar binnen gerichte resulterende
kracht zijn.

Die kracht is de gravitatiekracht.

We kunnen daarom al schrijven:

\[
\boxed{
F_g=ma_c.
}
\]

In een volgend hoofdstuk combineren we dit met de gravitatiewet:

\[
F_g=G\frac{Mm}{r^2}.
\]

Daaruit zullen onder andere de baansnelheid en omlooptijd van
satellieten volgen.

We stellen die volledige afleiding bewust uit: hier is het doel eerst
de wetten van Newton grondig te leren analyseren.


% ============================================================
% 22 DATA EN ICT
% ============================================================

\FVDsection{Onderzoek met data en ICT}

Newton II voorspelt bij constante massa:

\[
a=\frac{F_{\mathrm{res}}}{m}.
\]

Een experiment met een dynamiekwagentje kan dit verband onderzoeken.

Bij constante massa verwachten we:

\[
\boxed{
a\sim F_{\mathrm{res}}.
}
\]

Een grafiek van $a$ als functie van $F_{\mathrm{res}}$ hoort dus
lineair te zijn met helling:

\[
\frac1m.
\]

Bij constante kracht verwachten we:

\[
a\sim\frac1m.
\]

Een grafiek van $a$ tegen $m$ is dan niet lineair, maar een grafiek van
$a$ tegen $1/m$ wel.


\begin{oefening}[Newton II experimenteel onderzoeken]{\oefU}

Gebruik een dynamiekar, krachtensensor en bewegingssensor of een
geschikte simulatie.

\begin{question}
Formuleer een onderzoeksvraag over het verband tussen
$F_{\mathrm{res}}$, $m$ en $a$.
\end{question}

\begin{question}
Kies welke grootheid je varieert en welke je constant houdt.
\end{question}

\begin{question}
Voorspel de vorm van de grafiek vóór je meet.
\end{question}

\begin{question}
Verwerk de gegevens met een geschikte grafiek en trendlijn.
\end{question}

\begin{question}
Interpreteer de helling fysisch.
\end{question}

\begin{question}
Bespreek afwijkingen tussen model en experiment.
\end{question}

\end{oefening}


% ============================================================
% 23 GEINTEGREERDE OEFENINGEN
% ============================================================

\FVDsection{Geïntegreerde oefeningen}

\begin{oefening}[Van beweging naar kracht]{\oefU\ESS}

Een voorwerp met massa $3,0\,\mathrm{kg}$ beweegt langs de $x$-as met

\[
v_x(t)=4,0+2,5t.
\]

\begin{question}
Bepaal $a_x(t)$.
\end{question}

\begin{question}
Bepaal de resulterende kracht als functie van de tijd.
\end{question}

\begin{question}
Is de resulterende kracht constant? Motiveer.
\end{question}

\begin{question}
Welke soort rechtlijnige beweging voert het voorwerp uit?
\end{question}

\end{oefening}


\begin{oefening}[Van krachtfunctie naar beweging]{\oefG}

Een deeltje met constante massa $m$ beweegt in één dimensie. De
resulterende kracht wordt gegeven door:

\[
F_{\mathrm{res}}(t)=kt,
\]

met $k>0$.

Bij $t=0$ geldt:

\[
v(0)=v_0,
\qquad
x(0)=x_0.
\]

\begin{question}
Bepaal $a(t)$ met Newton II.
\end{question}

\begin{question}
Gebruik je kennis van integralen om $v(t)$ te bepalen.
\end{question}

\begin{question}
Bepaal vervolgens $x(t)$.
\end{question}

\begin{question}
Leg uit waarom deze beweging geen EVRB is.
\end{question}

\end{oefening}


\begin{oefening}[Robot op een vloer]{\oefG}

Een mobiele robot met massa $18\,\mathrm{kg}$ rijdt over een
horizontale vloer. Zijn aandrijving levert horizontaal een kracht van
$65\,\mathrm N$. De totale tegenwerkende kracht door rolweerstand en
andere verliezen wordt in een eerste model als constant
$20\,\mathrm N$ genomen.

\begin{question}
Maak een VLD.
\end{question}

\begin{question}
Bepaal de resulterende horizontale kracht.
\end{question}

\begin{question}
Bepaal de versnelling.
\end{question}

\begin{question}
De robot draagt daarna een extra lading van $12\,\mathrm{kg}$ terwijl
de twee krachten volgens het model gelijk blijven. Bepaal de nieuwe
versnelling.
\end{question}

\begin{question}
Beoordeel kritisch de aanname dat de tegenwerkende kracht bij de extra
lading onveranderd blijft.
\end{question}

\end{oefening}


\begin{oefening}[Foutenanalyse]{\oefU\ESS}

Een leerling lost een probleem op voor een kist die over een ruwe
horizontale vloer wordt getrokken en schrijft:

\[
F_N=mg,
\]

\[
F_w=\mu F_N,
\]

\[
F-F_w=ma.
\]

De trekkracht maakt in werkelijkheid een hoek van $30^\circ$ boven de
horizontale.

\begin{question}
Welke stap is problematisch?
\end{question}

\begin{question}
Maak het correcte VLD.
\end{question}

\begin{question}
Schrijf de correcte vergelijking in de verticale richting.
\end{question}

\begin{question}
Schrijf de correcte vergelijking in de horizontale richting.
\end{question}

\begin{question}
Leg uit waarom een goed krachtendiagram de fout vrijwel automatisch
zichtbaar maakt.
\end{question}

\end{oefening}


% ============================================================
% 24 VERDIEPING
% ============================================================

\FVDsection{Verdieping: Newton als differentiaalvergelijking}

Newton II kan voor een beweging langs één as geschreven worden als:

\[
F_{\mathrm{res}}(x,v,t)
=
m\frac{dv}{dt}.
\]

Omdat

\[
v=\frac{dx}{dt},
\]

kan ook:

\[
\boxed{
m\frac{d^2x}{dt^2}
=
F_{\mathrm{res}}(x,v,t).
}
\]

Dit is een differentiaalvergelijking: de onbekende is niet alleen een
getal, maar een functie $x(t)$.

Voor een constante resulterende kracht:

\[
F_{\mathrm{res}}=F_0
\]

krijgen we:

\[
\frac{d^2x}{dt^2}
=
\frac{F_0}{m}
=
\text{constant}.
\]

Na tweemaal integreren ontstaan precies de bekende
bewegingsvergelijkingen van de EVRB.

Zo zien we dat de kinematische formules geen losse regels zijn, maar
volgen uit Newton II wanneer de kracht constant is.


\begin{oefening}[Van Newton naar kinematica]{\oefG}

Neem:

\[
m\frac{d^2x}{dt^2}=F_0,
\]

met constante $F_0$.

\begin{question}
Toon dat

\[
a=\frac{F_0}{m}.
\]

\end{question}

\begin{question}
Integreer en gebruik $v(0)=v_0$ om te tonen dat

\[
v(t)=v_0+\frac{F_0}{m}t.
\]

\end{question}

\begin{question}
Integreer opnieuw en gebruik $x(0)=x_0$ om te tonen dat

\[
x(t)
=
x_0+v_0t+\frac12\frac{F_0}{m}t^2.
\]

\end{question}

\begin{question}
Leg uit hoe dit de verbinding tussen dynamica en de eerder bestudeerde
kinematica zichtbaar maakt.
\end{question}

\end{oefening}


% ============================================================
% 25 SAMENVATTING
% ============================================================

\FVDsection{Wat hebben we opgebouwd?}

De resulterende kracht is de vectoriële som van alle uitwendige
krachten op het gekozen systeem:

\[
\boxed{
\vec F_{\mathrm{res}}=\sum_i\vec F_i.
}
\]

Newton I:

\[
\boxed{
\sum\vec F=\vec 0
\Longrightarrow
\vec v=\text{constant}.
}
\]

Newton II:

\[
\boxed{
\sum\vec F=m\vec a.
}
\]

In componenten:

\[
\boxed{
\sum F_x=ma_x,
\qquad
\sum F_y=ma_y.
}
\]

Newton III:

\[
\boxed{
\vec F_{A\rightarrow B}
=
-\vec F_{B\rightarrow A}.
}
\]

Zwaartekracht nabij het aardoppervlak:

\[
\boxed{
F_g=mg.
}
\]

Statische wrijving:

\[
\boxed{
0\leq F_{w,s}\leq\mu_sF_N.
}
\]

Maximale statische wrijving:

\[
\boxed{
F_{w,s,\max}=\mu_sF_N.
}
\]

Kinetische wrijving:

\[
\boxed{
F_{w,k}=\mu_kF_N.
}
\]

Voor een cirkelbeweging:

\[
\boxed{
\sum F_r
=
ma_c
=
m\frac{v^2}{r}.
}
\]

De kern is echter niet het memoriseren van deze formules. De kern is
een situatie kunnen omzetten in een correct systeem, VLD en
vectorvergelijking.



% ============================================================
% 27 VOORUITBLIK
% ============================================================

\FVDsection{Vooruitblik}

We hebben nu twee krachtige ideeën samengebracht:

\[
\text{cirkelbeweging}
\qquad\text{en}\qquad
\text{Newton II}.
\]

Voor een cirkelbaan:

\[
\sum F_r=m\frac{v^2}{r}.
\]

Bij een satelliet kan de gravitatiekracht precies die radiale resultante
leveren.

In het volgende hoofdstuk keren we daarom terug naar gravitatie, maar
nu met nieuwe gereedschappen. We zullen onderzoeken hoe Newtons
gravitatiewet en de cirkelbeweging samen de beweging van manen,
planeten en satellieten kunnen verklaren en kwantificeren.

\end{document}
